\documentclass[main.tex]{subfiles}
\newcommand\chapterlabel{Ch-electrochemistry}\setcounter{figurenewcounter}{0}\setcounter{tablenewcounter}{0}\setcounter{formulanewcounter}{0}
\setcounter{figurenewcounter}{0}





\begin{document}
%\linenumbers
%\setcounter{chapter}{5}
  
  \import{files/}{ChapterName}


      \begin{marginfigure}
      \begin{tikzpicture} \node (a) at (0,0) {\includegraphics[width=4cm]{../Ch-electrochem/figure1}} node[rotate=90, font=\tiny] at ([yshift=.5cm,xshift=.1cm]a.south east) {\textsuperscript{\textcopyright} PxFuel} ;
\end{tikzpicture}
\end{marginfigure}


  \import{files/}{ChapterIntro}
\begin{marginfigure}%LEARNING GOALS BOX
\begin{mytcbox}{GOALS}
\begin{enumerate}[label=\protect\circled{\color{white}\arabic*}]
\item Identify anodes/cathodes
\item Calculate cell potentials
\item Interpret the line notation
\item Calculate cell potentials of concentration cells
\item Relate cell potential with $\Delta \text{G}^{\circ}$
\end{enumerate}
\end{mytcbox}
\vspace{1cm}
\begin{tcolorbox}[enhanced,colback=red!5!white,colframe=black!50!red,boxrule=1pt,
  arc=0pt,outer arc=0pt,drop heavy lifted shadow]
\faGears\ 
  \import{files/}{ChapterDiscussion}


\end{tcolorbox}
\end{marginfigure}%LEARNING GOALS BOX

\section{Introduction to galvanic cells}\import{files/}{SectionIntro-Introduction-to-galvanic-cells}
\sloppy\begin{description}
\item[\docfilehook{Components of a galvanic cell}{ }] 	\import{files/}{Subsection-Introduction-to-galvanic-cells-Components-of-a-galvanic-cell}
\item[\docfilehook{The electrodes: anode and cathode}{ }]\import{files/}{Subsection-Introduction-to-galvanic-cells-The-electrodes:-anode-and-cathode}
\item[\docfilehook{Cell potential}{ }]\import{files/}{Subsection-Introduction-to-galvanic-cells-Cell-potential}
\item[\docfilehook{Role of the salt bridge or the membrane}{ }]\import{files/}{Subsection-Introduction-to-galvanic-cells-Role-of-the-salt-bridge-or-the-membrane}
\item[\docfilehook{A galvanic cell example}{ }]\import{files/}{Subsection-Introduction-to-galvanic-cells-A-galvanic-cell-example}
\import{files/}{Figure-Daniell-cell}
\item[\docfilehook{The potentiometer}{ }] \import{files/}{Subsection-Introduction-to-galvanic-cells-The-potentiometer}
\import{problems/}{SampleProblem1}
\end{description}

\section{Standard reduction potentials}\import{files/}{SectionIntro-Standard-reduction-potentials} 
\sloppy\begin{description}
\item[\docfilehook{Electrode potentials}{ }]\import{files/}{Subsection-Standard-reduction-potentials-Electrode-potentials}
\item[\docfilehook{Standard conditions for reduction potentials}{ }]\import{files/}{Subsection-Standard-reduction-potentials-Standard-conditions-for-reduction-potentials}
\import{files/}{Figure-Batter-hydrogen-electrode}
\item[\docfilehook{Anodes and cathodes}{ }]\import{files/}{Subsection-Standard-reduction-potentials-Anodes-and-cathodes}
\import{problems/}{SampleProblem2}
\end{description}




\section{Line notation for galvanic cells}
There is a quick and easy way to represent a galvanic cell without having to draw the whole cell set up. This is called the line notation and galvanic cells are represented in a single line, starting from left to right. The anode is presented in the left, starting from the metallic part and followed by the electrolyte. A single line represents the liquid-solid contact. A double line represents the salt bridge and the cathode is represented in the right, starting for the electrolyte and finishing by the metal. As you can see, the line notation respect all interphase present in the cells: from left to right we have solid, liquid in contact with the salt bridge which is in contact with the liquid part of the cathode and finally we have the solid part of the cathode.
For example, the line notation of Daniell cell is:
\begin{center}\ce{Zn}| \ce{Zn^{2+}}\text{(1 M)}||\ce{Cu^{2+}}\text{(1 M)}|\ce{Cu}	\end{center}
In case there are several electrolytes in any of the electrodes, as all species are in liquid phase we separate them with just a comma. For example, in the galvanic cell below the cathode contains two different states of iron and uses Pt for the charge transfer:
\begin{center}\ce{Zn}| \ce{Zn^{2+}}\text{(1 M)}||\ce{Fe^{2+}}, \ce{Fe^{3+}}\text{(1 M)}|\ce{Pt}	\end{center}




\section{Standard reduction potential table}
The reduction potential of most of the electrodes are tabulated in standard conditions. This section will cover how to employ the table of standard potentials in order to obtain the electrode potential for a given electrode.
\sloppy
\begin{description}
\item[\docfilehook{Standard electrode potentials are expressed as reduction}{ }] 
Electrodes potentials can be expressed as reduction or as oxidation. This is because all electrodes can behave as anode or cathode, depending on the conterelectrode used in the galvanic cell. For example, for the \ce{Zn_(s)}/\ce{Zn^{2+}_{(aq)}} the reduction potential is:

\begin{center}\ce{Zn^{2+}_{(aq)} + 2e^- ->Zn_{(s)} }\hfill -0.76V\end{center}
whereas the oxidation potential is:
\begin{center}\ce{Zn_{(s)}		 ->	Zn^{2+}_{(aq)} + 2e^-		 }\hfill 0.76V\end{center}
In the table of standard potentials all reactions are expressed only as reduction and as such all potentials are reduction potentials. As such, all reaction will have electrons on the left side of the arrow. The electrode potential you will find on the table for the \ce{Zn_(s)}/\ce{Zn^{2+}_{(aq)}}  would be
\begin{center}\ce{Zn^{2+}_{(aq)} + 2e^- ->Zn_{(s)} }\hfill $\mathcal{E}^{\circ}=$-0.76V\end{center}


\import{files/}{Table-redox}









\end{description}
\end{document}
